

\SGARefTarget{sga3:iv:subsection:4:4}\subsection*{4.4. Exactness Properties of the Category of Sheaves}
\addcontentsline{toc}{subsection}{\SGARefLink{sga3:iv:subsection:4:4}{4.4}. Exactness Properties of the Category of Sheaves}

\medskip
\noindent\SGARefTarget{sga3:iv:theorem:4:4:1}\textbf{Theorem 4.4.1.}
\begin{enumerate}
  \item[(i)]\SGARefTarget{sga3:iv:theorem:4:4:1:item:i:01} \textit{Arbitrary inverse limits exist in
    \(\widetilde{\mathscr C}\), and \textquotedblleft are computed in
    \(\widehat{\mathscr C}\)\textquotedblright; that is, the inclusion functor
    \[
      i:\widetilde{\mathscr C}\longrightarrow\widehat{\mathscr C}
    \]
    commutes with inverse limits. More explicitly, if \((X_\alpha)\) is an
    inverse system of sheaves, the presheaf}
    \[
      \varprojlim_\alpha i(X_\alpha):
      S\longmapsto\varprojlim_\alpha X_\alpha(S)
    \]
    \textit{is a sheaf, and}
    \[
      i\!\left(\varprojlim_\alpha X_\alpha\right)
      =\varprojlim_\alpha i(X_\alpha).
    \]

  \item[(ii)]\SGARefTarget{sga3:iv:theorem:4:4:1:item:ii:01} \textit{Arbitrary direct limits exist in
    \(\widetilde{\mathscr C}\). If \((X_\alpha)\) is a direct system of
    sheaves, then}
    \[
      \varinjlim_\alpha X_\alpha
      =a\!\left(\varinjlim_\alpha i(X_\alpha)\right),
    \]
    \textit{where \(\varinjlim_\alpha i(X_\alpha)\) is the presheaf direct
    limit of the \(i(X_\alpha)\):}
    \[
      \varinjlim_\alpha i(X_\alpha):
      S\longmapsto\varinjlim_\alpha X_\alpha(S).
    \]

  \item[(iii)]\SGARefTarget{sga3:iv:theorem:4:4:1:item:iii:01} \textit{The functor}
    \[
      a:\widehat{\mathscr C}\longrightarrow\widetilde{\mathscr C}
    \]
    \textit{commutes with arbitrary direct limits and finite inverse limits.}
\end{enumerate}

Assertions \SGARefLink{sga3:iv:theorem:4:4:1:item:i:01}{(i)} and \SGARefLink{sga3:iv:theorem:4:4:1:item:ii:01}{(ii)} follow formally from the adjunction formula
(\SGARefLink{sga3:iv:theorem:4:3:14}{4.3.14}), and\footnote{Editor's note: the original has been modified here.}
the first assertion of \SGARefLink{sga3:iv:theorem:4:4:1:item:iii:01}{(iii)} follows from \SGARefLink{sga3:iv:theorem:4:4:1:item:ii:01}{(ii)}. The second assertion of \SGARefLink{sga3:iv:theorem:4:4:1:item:iii:01}{(iii)}
was already noted in \SGARefLink{sga3:iv:remark:4:3:16}{4.3.16}.

\medskip
\noindent\SGARefTarget{sga3:iv:scholium:4:4:2}\textbf{Scholium 4.4.2.}
\textit{\SGARefLink{sga3:iv:theorem:4:4:1}{This theorem} makes it possible to use the following method for proving
in \(\widetilde{\mathscr C}\) a statement involving both arbitrary direct
limits and finite inverse limits (for example,
\textquotedblleft every epimorphism is universally
effective\textquotedblright; cf. below).
One begins by proving the corresponding statement in the category of sets,
then extends it pointwise to the category of presheaves. One then uses \SGARefLink{sga3:iv:theorem:4:4:1}{the
preceding theorem} to pass from the category of presheaves to the category of
sheaves. Many examples of this method will occur below (\SGARefLink{sga3:iv:proposition:4:4:3}{4.4.3}, \SGARefLink{sga3:iv:lemma:4:4:6}{4.4.6}, \SGARefLink{sga3:iv:proposition:4:4:9}{4.4.9},
and so forth).}

Finally, observe that statements concerning the category of presheaves are
formally corollaries of the corresponding statements concerning the category
of sheaves. It is enough to take the coarsest
(\textquotedblleft chaotic\textquotedblright)
topology, namely the topology defined by
\[
  J(S)=\{S\}
  \qquad\text{for every }S\in\operatorname{Ob}\mathscr C;
\]
every functor is a sheaf for this topology.

\medskip
\noindent\SGARefTarget{sga3:iv:proposition:4:4:3}\textbf{Proposition 4.4.3.}
\textit{Let \(\mathscr F=\{F_i\to F\}\) be a family of morphisms of sheaves.
The following conditions are equivalent:}
\begin{enumerate}
  \item[(i)]\SGARefTarget{sga3:iv:proposition:4:4:3:item:i:01} \textit{\(\mathscr F\) is an epimorphic family;}
  \item[(ii)]\SGARefTarget{sga3:iv:proposition:4:4:3:item:ii:01} \textit{\(\mathscr F\) is a universal effective epimorphic
    family (\SGARefLink{sga3:iv:remark:1:13}{1.13});}
  \item[(iii)]\SGARefTarget{sga3:iv:proposition:4:4:3:item:iii:01} \textit{\(\mathscr F\) is covering (\SGARefLink{sga3:iv:definition:4:2:3}{4.2.3});}
  \item[(iv)]\SGARefTarget{sga3:iv:proposition:4:4:3:item:iv:01} \textit{the sheaf image of \(\mathscr F\), that is, the sheaf
    associated with the presheaf image of \(\mathscr F\) (\SGARefLink{sga3:iv:definition:4:1:2}{4.1.2}), is \(F\).}
\end{enumerate}

The equivalence of \SGARefLink{sga3:iv:proposition:4:4:3:item:iii:01}{(iii)} and \SGARefLink{sga3:iv:proposition:4:4:3:item:iv:01}{(iv)} follows from \SGARefLink{sga3:iv:corollary:4:3:12}{4.3.12}. The remaining
equivalences will follow from the next lemmas.

\medskip
\noindent\SGARefTarget{sga3:iv:lemma:4:4:4}\textbf{Lemma 4.4.4.}
\textit{Let \(f:X\to Y\) be a monomorphism of sheaves that is also an
epimorphism. Then \(f\) is an isomorphism.}

\SGARefLink{sga3:iv:lemma:4:4:4}{The lemma} is clear first of all in the category of sets. Let us next prove it
in the category of presheaves. Consider the presheaf
\[
  V:S\longmapsto Y(S)\amalg_{X(S)}Y(S);
\]
this is the amalgamated sum of \(Y\) and \(Y\) over \(X\) in the category of
presheaves.\footnote{Editor's note: the remainder of the proof has been
slightly modified.}
Let \(i_1,i_2:Y\rightrightarrows V\) be the two coordinate morphisms. If
\(X\to Y\) is an epimorphism in the category of presheaves, then \(i_1=i_2\).
In that case, for every \(S\), the map \(X(S)\to Y(S)\) is surjective; since
it is also injective, it is bijective, and hence \(X\to Y\) is an
isomorphism.

Finally, work in the category \(\widetilde{\mathscr C}\) of sheaves. By
\SGARefLink{sga3:iv:theorem:4:4:1:item:ii:01}{4.4.1(ii)}, the amalgamated sum in \(\widetilde{\mathscr C}\) of the sheaves
\(Y\) and \(Y\) over \(X\) is the sheaf \(Z\) associated with the presheaf
\(V\). Consider the diagram
\[
  X\xrightarrow{f}Y
  \overset{i_1}{\underset{i_2}{\rightrightarrows}}
  V\xrightarrow{\tau}Z=a(V).
\]
We have \(i_1f=i_2f\), hence \(\tau i_1f=\tau i_2f\), and therefore
\(\tau i_1=\tau i_2\), since \(f\) is an epimorphism in
\(\widetilde{\mathscr C}\). By part \SGARefLink{sga3:iv:lemma:4:4:5:item:iii:01}{(iii)} of \SGARefLink{sga3:iv:lemma:4:4:5}{the following lemma}, the
presheaf \(V\) is separated, so \(\tau\) is a monomorphism (\SGARefLink{sga3:iv:proposition:4:3:11:item:iv:01}{4.3.11(iv)}).
Thus \(i_1=i_2\), and we saw above that this implies that \(f:X\to Y\) is an
isomorphism. This proves Lemma~\SGARefLink{sga3:iv:lemma:4:4:4}{4.4.4}, once it has been verified that \(V\) is
separated.
